\chapter{Implemented cryptography and verifiable artifact contracts}
\label{ch:implemented-cryptography}

\reviewstatus{Executable library-backed encryption, signature and key-encapsulation
interfaces, with 25 passing regression-test methods and preserved runtime evidence.
The framing and tests are local contributions. The cryptographic primitives are
external implementations; no new security reduction or Lean verification of those
implementations is claimed.}

\section{From mathematical examples to executable interfaces}

The finite curve, hidden-relation and noise calculations in
Chapters~\ref{ch:quantum-crypto} and~\ref{ch:pqc} explain mathematical mechanisms.
The new \path{extension/crypto/} directory supplies a separate operational layer:
authenticated artifact encryption, signatures and post-quantum key encapsulation.
Its Python interface delegates every cryptographic primitive to
\texttt{cryptography}~50.0.0 and its native backend. This choice makes the
implementation boundary inspectable without presenting an educational arithmetic
example as a deployed cryptosystem \cite{CryptoAPI50AEAD,CryptoAPI50Ed25519,CryptoAPI50MLKEM}.

\begin{longtable}{p{.23\textwidth}p{.28\textwidth}p{.39\textwidth}}
\toprule
Function family & Primitive & Contract exercised\\\midrule
Artifact encryption & AES-256-GCM; ChaCha20-Poly1305 & Byte confidentiality and
authentication of ciphertext plus public version, algorithm and context\\
Artifact signatures & Ed25519 & Verification of context-framed bytes relative
to a separately trusted public key\\
Key encapsulation & ML-KEM-768; ML-KEM-1024 & Creation and recovery of a shared
secret through the backend's standard encapsulation API\\\bottomrule
\end{longtable}

The algorithm sources are the GCM recommendation, the ChaCha20--Poly1305 and
EdDSA RFCs, and FIPS~203 for ML-KEM
\cite{CryptoGCM38D,CryptoRFC8439,CryptoRFC8032,CryptoFIPS203}.
ML-KEM's module-lattice basis connects conceptually to the preceding order-lattice
discussion. Its security is not derived here from a finite determinant, a Hodge
gap, a tropical model, GRH, or the fluid construction. A standard algorithm and a
FIPS-validated cryptographic module are also different claims: this release
records no module-validation certificate.

\section{Authenticated serialization and expected context}

The encryption envelope contains exactly five JSON fields:
\texttt{version}, \texttt{algorithm}, \texttt{context}, \texttt{nonce} and
\texttt{ciphertext}. The parser rejects unknown or duplicate fields, incorrect
types, unsupported versions, noncanonical base64 and invalid lengths. Version is
the integer~1; even the JSON Boolean \texttt{true} is rejected instead of being
accepted through Python's integer comparison. The plaintext limit is 16~MiB and
the context contains 1--1024 UTF-8 bytes. Context is authenticated public metadata,
so confidential information belongs in the encrypted payload.

Associated data consists of a fixed application-domain prefix and the compact,
sorted-key JSON encoding of version, algorithm and context. The recipient supplies
the expected algorithm and expected context independently of the incoming
envelope. These expectations cannot safely be copied from attacker-controlled
metadata. Altering a header and simultaneously selecting the altered expectation
still fails the library's authentication check. The full 16-byte authentication
tag is retained, and no plaintext is returned before verification succeeds.
These are concrete wrapper contracts around the external AEAD APIs, not an
independent proof of those APIs' computational security \cite{CryptoAPI50AEAD}.

Each encryption generates a fresh 96-bit random nonce internally. The interface
does not accept caller-selected nonces. For $m$ ideal independent uniform draws,
the union bound over pairs gives
\[
 \Pr[\text{a nonce collision}]\leq
 \binom m2\,2^{-96}=\frac{m(m-1)}{2^{97}}.
\]
For $m\leq2^{20}$ this particular probability is below $2^{-57}$. This calculation
addresses nonce collisions only; it is not the complete AEAD forgery or privacy
bound. Random sampling does not guarantee perpetual uniqueness. The stateless
wrapper cannot enforce a distributed key-use quota or retain a counter across
process restarts. Applications need their own key lifecycle and usage policy,
with independent keys for different algorithms and purposes. A fresh key for a
single artifact is the simplest example, while secure storage and distribution
of that key remain application responsibilities.

\section{Signatures, replay and scientific meaning}

The signature wrapper signs a separate domain prefix, the four-byte big-endian
length of the UTF-8 context, the context itself, and the exact artifact bytes.
The explicit length prevents the ambiguity between, for example, context
\texttt{a} followed by message \texttt{bc} and context \texttt{ab} followed by
message \texttt{c}. This is ordinary Ed25519 applied to the framed byte string;
it is not the distinct Ed25519ctx variant specified in the EdDSA family
\cite{CryptoRFC8032,CryptoAPI50Ed25519}.

Verification requires an authenticated association between the public key and
the intended signer. A valid signature authenticates bytes relative to that key;
it does not establish that the signed mathematical assertion is true. Signing
a verification manifest can preserve provenance, while its actual proof or
test evidence must still be checked. The example keys are ephemeral and do not
constitute an author identity or an archival trust anchor.

An authentic old envelope or signature remains authentic when replayed.
Putting an artifact identifier or release number in the expected context binds
that value, but preventing a second acceptance requires receiver state.
Likewise, byte-exact signatures do not automatically identify semantically
equivalent JSON encodings; any required canonicalization is an application
contract. No secret key belongs in a public repository or client-side static
site, and Python object destruction is not a guarantee of memory zeroization.

\section{The post-quantum boundary}

The two ML-KEM interfaces return and recover 32-byte shared secrets using the
native library. The exercised ML-KEM-768 public-key and ciphertext lengths are
1184 and 1088 bytes; ML-KEM-1024 uses 1568 bytes for each. The tests serialize
and reload public keys, encapsulate, decapsulate and check these sizes against
the documented API \cite{CryptoAPI50MLKEM}.

ML-KEM is a key-encapsulation mechanism, not a signature scheme or complete
authenticated transport protocol \cite{CryptoFIPS203}. It does not establish
either party's identity, specify file encryption or prevent replay. In
particular, an invalid ciphertext of the correct length normally yields a
rejection secret through implicit rejection; a returned byte string is not
evidence that the ciphertext is valid. The regression checks this behavior and
checks that malformed lengths are rejected. Combining these functions into a
KEM/DEM transport would require separately specified authentication, transcript
binding, derivation and lifecycle rules; this release makes no such protocol
claim. Ed25519 remains a classical signature, so its presence beside ML-KEM does
not make every component post-quantum.

\section{Recorded execution and formal coverage}

The recorded execution used Python~3.12.14,
\texttt{cryptography}~50.0.0 and OpenSSL~4.0.1. All five algorithm choices in the
table were exercised. The 25 test methods passed with no failures, errors or
skips; parameterized cases additionally cover empty, binary, Unicode and larger
payloads; ciphertext, tag and nonce changes; wrong keys, contexts and algorithms;
parser ambiguity; signature changes; context framing; both KEM parameter sets;
implicit rejection; and incorrect KEM lengths. Private keys and shared secrets
were generated in memory and were neither persisted nor printed.

\begin{lstlisting}
python extension/crypto/verify_crypto.py
\end{lstlisting}

The reproducible evidence is \path{extension/crypto/verification.json} together
with \path{extension/crypto/verification.log}. The JSON record identifies the
runtime versions and hashes the exact implementation, test driver, dependency
pin and test log. A successful finite test run is not a proof for all inputs,
a security reduction, an independent backend audit or a certification claim.

The companion Lean noise-cancellation and hidden-relation lemmas concern the
finite algebra stated in their types. They do not formalize this Python code,
the native backend, signature unforgeability, chosen-ciphertext security,
side-channel resistance or authenticated key distribution. Keeping those
interfaces explicit permits useful connections between arithmetic research and
cryptographic engineering without changing the evidentiary meaning of either.
